%%%%%%%%%%%%%%%%%%%%%%%%%%%%%%%%%%%%%%%%
%% HiMCM/MidMCM LaTeX Template -- 2026 %%
%% Problem B: Requirements 1/2 and partial Requirement 4 working draft
%% Efficiency/effectiveness scope separation, November 8, 2026.
%%%%%%%%%%%%%%%%%%%%%%%%%%%%%%%%%%%%%%%%
\documentclass[12pt]{article}
\usepackage{geometry}
\geometry{left=1in,right=0.75in,top=1in,bottom=1in}
\newcommand{\Problem}{B}
\newcommand{\Team}{1111111} % Replace with the actual Team Control Number.
\usepackage{newtxtext}
\usepackage{amsmath,amssymb,amsthm}
\usepackage{newtxmath}
\usepackage[pdftex]{graphicx}
\usepackage{xcolor}
\usepackage{fancyhdr}
\usepackage{microtype}
\usepackage{xurl}
\usepackage{needspace}
\usepackage{booktabs,tabularx,array}
\usepackage[hidelinks]{hyperref}
\setlength{\headheight}{15pt}
\setlength{\parskip}{3pt}
\setlength{\parindent}{1.2em}
\widowpenalty=10000
\clubpenalty=10000
\lhead{Team \Team}
\rhead{}
\cfoot{}
\newtheorem{theorem}{Theorem}
\newtheorem{corollary}[theorem]{Corollary}
\newtheorem{lemma}[theorem]{Lemma}
\newtheorem{definition}{Definition}

\begin{document}
\graphicspath{{.}}
\DeclareGraphicsExtensions{.pdf,.jpg,.png}
\thispagestyle{empty}
\vspace*{-16ex}
\centerline{\begin{tabular}{*3{c}}
\parbox[t]{0.3\linewidth}{\begin{center}\textbf{Problem Chosen}\\ \Large \Problem\end{center}}&
\parbox[t]{0.3\linewidth}{\begin{center}\textbf{2026\\ HiMCM/MidMCM\\ Summary Sheet}\end{center}}&
\parbox[t]{0.3\linewidth}{\begin{center}\textbf{Team Control Number}\\ \Large \textcolor{red}{\Team}\end{center}}\\
\hline
\end{tabular}}

\begin{center}
\vspace{1.5em}
{\Large\bfseries Measuring Useful Cooling in Open Systems}\\[0.8em]
{\bfseries Working summary: efficiency and its effectiveness extension}
\end{center}
An open freezer case must maintain appropriate product temperatures despite heat exchange with its surroundings. Electricity use alone does not show how well it performs this task.

We distinguish thermal demand, electrical input, and delivered storage service. Infiltration, cabinet heat transfer, internal equipment, restocking, and defrost contribute to demand. Equipment condition modifies performance, while temperature differences across product locations determine whether the intended service is delivered.

We identify four quantities that define cooling service efficiency: energy input, service quantity, local service quality, and service duration. Compliant product mass-hours per kilowatt-hour provide one possible freezer measure. Product-specific temperature requirements remain a separate adequacy check. Comparisons require prescribed loading, matched ambient conditions, and complete operating cycles.

We extend this principle to a regular-size New York City bus shelter, using documented dimensions and a three-zone model. Fan-driven air movement and evaporative temperature/humidity changes feed into a common SET-based hot-discomfort proxy. Signed relief accumulated over passenger-hours is divided by electrical energy. SET defines the useful-cooling numerator; it is not an additional optimization objective in the efficiency comparison.

The bus model is the fixed-background limit of a shared heat/moisture and local-delivery framework. In the illustrative dry and humid cases, A has the highest service efficiency. A delivery-loss stress test identifies conditions where B becomes more efficient. The power ratios give transparent break-even requirements: B and C must deliver respectively 4.8 and 10 times A's accumulated relief to exceed its efficiency. These are conditional comparisons, not measured equipment rankings.

For Canberra MRT station, the same device, delivery and SET modules are retained while background air and surfaces evolve with ventilation and passenger loads. Coverage weights the people who receive each local condition. Among all 24 sampled installations, A3 has the highest efficiency, $37.58$; among six benchmarks with equal $68.2\%$ direct coverage, A36 is highest at $35.49$. These are distinct comparison sets. Local-delivery and heat-allocation sensitivities test the efficiency ranking.

Requirement 4 separately addresses overall effectiveness. A3's $5.68\%$ direct coverage illustrates why the efficiency leader need not provide adequate service. We relocate residual-discomfort and energy tradeoffs to that extension and formulate visitor-experience and cost inputs there. The cost and non-thermal experience inputs are not yet calibrated, so no overall-effectiveness winner is claimed.

\medskip
\noindent\textit{Draft status.} Requirements 1 and 2 are developed; Requirement 4 has a partial decision framework. Requirement 3 cross-venue applications, the cost/experience calibration for Requirement 4, Requirement 5 and the final summary remain unfinished. Numerical results are simulations under explicit assumptions.

\clearpage
\pagestyle{fancy}
\rhead{Page \thepage}
\tableofcontents
\clearpage

\section{Understanding Cooling Efficiency}
\label{sec:understand}

An open freezer case succeeds by preserving its products. Low electricity use may reflect a smaller workload, shorter operation, or inadequate temperatures; high heat removal may reflect unwanted environmental heat gains. Cooling efficiency must therefore relate energy input to the storage service actually delivered.


Figure~\ref{fig:open-case} makes the open-front geometry concrete: the front opening permits air exchange, the cabinet surfaces transmit heat, and the shelves separate product locations whose temperatures should be checked. The pictured model is a medium-temperature refrigerated display case \cite{hillphoenix}; its geometry illustrates these mechanisms, while freezer operating temperatures and product requirements must be specified separately.

\begin{figure}[!ht]
\centering
\includegraphics[width=0.76\linewidth]{figures/09_open_display_case.pdf}
\caption{Open-front refrigerated display case: user-supplied schematic with external dimensions matching the Hillphoenix CF-R 4-ft configuration \cite{hillphoenix}. Millimetre labels are rounded; the internal shelf arrangement is schematic.}
\label{fig:open-case}
\end{figure}

\subsection{Selecting the quantities to evaluate}

Table~\ref{tab:service} identifies four quantities that enter the useful-service numerator or energy denominator in both settings.

\begin{table}[!ht]
\centering
\caption{Evaluation quantities and their transfer to a bus stop.}
\label{tab:service}
\small
\renewcommand{\arraystretch}{1.18}
\begin{tabularx}{\linewidth}{@{}>{\raggedright\arraybackslash}p{0.19\linewidth}
>{\raggedright\arraybackslash}X >{\raggedright\arraybackslash}X@{}}
\toprule
Quantity & Open freezer case & Bus-stop shelter\\
\midrule
Energy input & Electricity used by the case and its attributable refrigeration plant (kWh). & Electricity used by fans, pumps, and other cooling equipment (kWh).\\
Service quantity & Actual mass of products stored under a specified loading protocol (kg). & Number of people present in each occupied zone.\\
Local service quality & Product temperatures relative to product-specific storage requirements. & Thermal discomfort relative to the same setting without active cooling.\\
Service duration & Duration of product storage during the evaluation period (h). & Duration of passenger exposure while waiting (h).\\
\bottomrule
\end{tabularx}
\end{table}

\noindent\textbf{Energy input} includes auxiliaries and attributable plant electricity without overlapping meter readings. \textbf{Service quantity} accounts for workload, subject to prescribed loading limits.

\textbf{Local service quality} requires product temperatures at representative locations, rather than discharge-air temperature alone \cite{ashrae}. Averages can hide a warm shelf, and uniformity does not guarantee acceptable storage; product-temperature requirements remain essential \cite{etl}.

\textbf{Service duration} distinguishes sustained performance from a favorable snapshot. Include restocking, compressor cycling, and defrost recovery, with comparable initial and final thermal states so that stored cold cannot inflate the result.


\subsection{Factors that explain performance}

The explanation follows the path from environmental loads to local conditions and then to delivered service. Its process-oriented organization draws on a supplied AI analysis, with the physical dependencies revised here \cite{gemini}.

\noindent\textbf{From surroundings to local conditions.} Infiltration carries sensible heat and moisture through the opening; walls transmit heat, and exposed surfaces exchange radiation. Geometry, insulation, ambient temperature, humidity, and airflow therefore alter the conditions around the products and the refrigeration demand \cite{ashrae}.

\noindent\textbf{From operation to local conditions.} Lights and fans release heat, warmer incoming stock adds a temporary load, and defrost interrupts normal cooling \cite{ashrae}. Equipment condition changes insulation, airflow, and power demand. These mechanisms affect both local conditions and measured electricity; age itself is not an additional score.

\noindent\textbf{From local conditions to service.} Product placement and thermal response determine whether the food remains within its required range. Accumulating compliant product mass over time gives useful storage service. The same sequence applies to passengers, but their response depends on the thermal environment and their exposure, rather than food-temperature limits.

\subsection{Implications for the bus-stop model}

The preceding analysis motivates a common evaluation principle:
\begin{equation}
\boxed{\text{Cooling efficiency}
=\frac{\text{useful service accumulated over objects and time}}
{\text{energy input over the same period}}.}
\label{eq:principle}
\end{equation}

For a prescribed product load, one concrete freezer measure is
\begin{equation}
 \eta_f=\frac{\int\sum_i m_i(t)\mathbf{1}\{T_{p,i}(t)\in\mathcal T_i\}\,dt}
 {\int p_f(t)\,dt},
\label{eq:freezer-service}
\end{equation}
where $\mathcal T_i$ is the product-specific admissible temperature range, $t$ is in hours and $p_f$ is in kW. Power includes attributable plant and auxiliary loads; comparable loading and complete operating cycles prevent a smaller workload or stored cold from inflating the score. Temperature excursions remain a separate adequacy check. Human service instead measures discomfort relief over passenger-hours. The two numerators describe different recipients and cannot be numerically ranked against one another.

For passengers, a fan changes air movement while evaporative equipment changes temperature and humidity \cite{energystar,evap}. Weight the response by where people stand and how long they wait; empty space earns no service credit. In Requirements 2 and 3, compare $\eta=U/E$ under matched weather, loading and evaluation periods. Partial relief enters $U$; interpreting whether that relief is adequate, and balancing it with visitor experience and cost, belongs to the separate effectiveness extension in Section~\ref{sec:effectiveness}.

\subsection{A common physical structure and a controlled extension}

All three settings exchange heat and moisture with their surroundings. For a fixed dry-air mass $M_z$ and balanced incoming/outgoing air streams, a common control-volume structure is
\begin{align}
 M_z\dot h_z&=\sum_j\dot m_{jz}(h_j-h_z)+S_{h,z},\\
 M_z\dot w_z&=\sum_j\dot m_{jz}(w_j-w_z)+S_{w,z}.
\label{eq:common-balance}
\end{align}
Here $h$ is J/kg dry air and $w$ is kg vapor/kg dry air; $S_h$ and $S_w$ collect the consistently defined heat and water sources. Refrigeration, product heat transfer and possible condensation enter a freezer's sources. Fans, evaporation, people and surfaces enter the public-space sources. The recipient-response module changes from food-temperature compliance to human SET; conservation and service accounting remain the organizing principles. Requirement 1 establishes these boundaries rather than claiming a calibrated freezer simulation.

\begin{figure}[!ht]
\centering
\includegraphics[width=\linewidth]{figures/19_model_inheritance.pdf}
\caption{Model inheritance: common conservation and service accounting, object-specific responses, and a controlled extension from fixed to dynamic background air.}
\label{fig:inheritance}
\end{figure}

The bus model treats the outdoor air as a fixed reservoir and resolves steady local delivery. The rail model retains that local calculation but solves the background balances and surface storage over time. This follows the heat/moisture balance structure used in building simulation \cite{epzone}. It also takes a process-level lesson from Cruz-Salas et al.\ \cite{cruz2026}: distinguish evaporative treatment, air transport and the conditions in the occupied space. Their study concerns misting in a dry-season, unoccupied Mexican cafeteria; its water-loss factors and measured cooling are not parameters for our wet-pad devices or Singapore case.

\Needspace{9\baselineskip}
\section{Cooling a Bus Stop: From Local Conditions to Passenger Service}
\label{sec:bus}

We apply the service principle in Section~\ref{sec:understand} to Requirement 2(a). A device first changes the environment at each passenger's position; those changes alter thermal discomfort; the resulting relief is accumulated over people and time. This separates three questions: how much relief is delivered, where it is delivered, and how much electricity it requires.


We adopt the regular-size New York City shelter in Figure~\ref{fig:shelter-geometry} as the common geometry for all three devices. Its documented dimensions provide a reproducible physical scale for the occupied zones.

\begin{figure}[!ht]
\centering
\includegraphics[width=0.82\linewidth]{figures/07_bus_shelter_geometry.pdf}
\caption{Adopted baseline geometry: regular-size New York City bus shelter. User-supplied schematic; dimensions verified against NYC DOT's \emph{Street Design Manual} \cite{nycshelter}. Millimetre labels are rounded.}
\label{fig:shelter-geometry}
\end{figure}

The official dimensions are $14$ ft by $5$ ft by $8$ ft $11$ in \cite{nycshelter}. Converting the imperial values before rounding gives $L=4.2672$, $W=1.5240$, and $H=2.7178$ m. We use the gross plan area $A=LW=6.5032\,\mathrm{m^2}$ and divide the length into three equal strips, each with $A_z=A/3=2.1677\,\mathrm{m^2}$. The bench and structure occupy part of this footprint; $A$ is a geometric reference, not a net standing area or a capacity estimate.

\begin{figure}[!ht]
\centering
\includegraphics[width=0.82\linewidth]{figures/08_shelter_zoning.pdf}
\caption{Calculation zones derived from the adopted dimensions. Equal strips and the delivery reference at $x=0$ are modeling choices; dots mark representative zone coordinates, not passenger locations. No bench or obstruction dimensions are inferred.}
\label{fig:shelter-zones}
\end{figure}

Representative coordinates are $(x_z,W/2)$, with $x_z=(z-\tfrac12)L/3$. Delivery is assumed to originate at $x=0$. Geometry defines the accounting regions; it does not determine the assumed speeds and mixing fractions.

\subsection{Data, interpretation, and a reproducible scenario}

The problem gives apparent cooling up to $3\,^{\circ}\mathrm C$ for A and air-temperature reductions up to $5/8\,^{\circ}\mathrm C$ for B/C \cite{problem}. We do not equate these quantities or treat A's unspecified ``feels like'' effect as measured SET. We interpret the table's kW/hr entries as mean powers $p_A=0.25$, $p_B=1.2$, $p_C=2.5$ kW; a one-hour comparison uses $0.25/1.2/2.5$ kWh. This is a unit interpretation, not verified equipment specifications.

Airflow, coverage, rating weather and evaporative effectiveness are unspecified. Tables~\ref{tab:bus-assumptions}--\ref{tab:bus-fields} define a reproducible scenario rather than identify actual device performance.

\begin{table}[!ht]
\centering\small
\caption{Common inputs: documented geometry and explicit scenario assumptions.}
\label{tab:bus-assumptions}
\renewcommand{\arraystretch}{1.12}
\begin{tabularx}{\linewidth}{@{}p{0.31\linewidth}X@{}}
\toprule
Quantity & Value, basis, and interpretation\\
\midrule
Shelter & NYC regular-size geometry: $4.2672\times1.5240\times2.7178$ m \cite{nycshelter}; three equal $2.1677\,\mathrm{m^2}$ gross zones. Full shade is assumed.\\
Weather & $T_a=34\,^{\circ}\mathrm C$; dry/humid cases RH $=30/70\%$; sweep $20$--$80\%$.\\
Radiation and background air & Mean radiant temperature $T_r=36\,^{\circ}\mathrm C$; baseline local speed $v_0=0.20$ m/s; pressure $101325$ Pa.\\
Reference passenger & Standing, $1.2$ met, $0.5$ clo; default body area $1.8258\,\mathrm{m^2}$ and zero external work in the SET implementation.\\
Occupancy and period & Assumed mean simultaneous occupancy $(4,4,4)$ over one hour: $12$ person-hours; an illustrative load, not a measured count or rated capacity.\\
Alternative distributions & Concentration stress cases $(8,3,1)$ and $(1,3,8)$; both retain the same total of 12 people.\\
Service reference & SET threshold $\Theta_c=26\,^{\circ}\mathrm C$, an assumed hot-discomfort reference; tested at $24$ and $28\,^{\circ}\mathrm C$.\\
\bottomrule
\end{tabularx}
\end{table}

\begin{table}[!ht]
\centering\small
\caption{Assumed local delivery profiles. Entries follow near/middle/edge order.}
\label{tab:bus-fields}
\begin{tabularx}{\linewidth}{@{}lXXc@{}}
\toprule
Device & Local air speed (m/s) & Treated-stream fraction $\beta_{s,z}$ & $\varepsilon_s$\\
\midrule
A & $1.00 / 0.55 / 0.25$ & immaterial at $\chi=0$ & ---\\
B & $0.80 / 0.50 / 0.30$ & $0.80 / 0.45 / 0.10$ & $0.70$\\
C & $1.00 / 0.80 / 0.50$ & $0.90 / 0.70 / 0.40$ & $0.70$\\
\bottomrule
\end{tabularx}
\end{table}

The entries represent zone averages over the strips in Figure~\ref{fig:shelter-zones}; they are not measured NYC delivery data. Here $\beta_{s,z}$ is a local dry-air mixing fraction, not the share of total device flow assigned to a zone; its entries need not sum to one. The larger coverage assumed for C is tested, not presented as a fact supplied by the problem.

\subsection{Predicting the local environment}

\noindent\textbf{Fan A: air movement.} We set $T_{A,z}=T_a$ and retain ambient moisture, while prescribing local speeds from Table~\ref{tab:bus-fields}. The resulting convection and evaporation enter the passenger-response model. This follows the physical distinction between moving air and lowering dry-bulb temperature \cite{fanguide}. The introductory bus calculation separately neglects local motor/pump heating. A fixed outdoor background alone does not justify that approximation; we relax it in the heat-allocation sensitivity and in the rail extension.

\noindent\textbf{B and C: evaporation followed by dilution.} Direct evaporation reduces dry-bulb temperature while increasing moisture; the entering wet-bulb temperature $T_{wb}$ limits the achievable depression \cite{evapequip}. We adopt
\begin{equation}
 T_{o,s}=T_b-\min\{\Delta T_s^{\max},\;\varepsilon_s(T_b-T_{wb,b})\},
 \qquad s\in\{B,C\},
\label{eq:outlet}
\end{equation}
where the inlet/background state $b$ is outdoor air for the bus and evolving local background air for the rail station. Here $\Delta T_B^{\max}=5$ and $\Delta T_C^{\max}=8\,^{\circ}\mathrm C$. Applying these advertised shelter-temperature reductions as outlet caps is an explicit modeling convention; the problem does not specify their measurement location. The caps do not imply that either device achieves them under every climate.

Let $w$ denote kg of water vapor per kg of dry air and $h$ the moist-air enthalpy in J/kg of dry air. Using standard SI psychrometric relations \cite{psychrolib},
\begin{align}
 w&=0.621945\frac{p_v}{p-p_v},\qquad p_v=\phi p_{sat}(T),\\
 h(T,w)&=1006T+w(2501000+1860T),
\label{eq:enthalpy}
\end{align}
where $\phi$ is RH as a fraction. Approximating evaporation as isenthalpic, $h_{o,s}\simeq h_b$, gives
\begin{equation}
 w_{o,s}=\frac{h_b-1006T_{o,s}}{2501000+1860T_{o,s}}.
\end{equation}
This neglects fan heat and make-up-water enthalpy; near-constant enthalpy is not an assertion of exact constant wet-bulb temperature \cite{evap}.

\noindent\textbf{One local-delivery model for both venues.} Let a device stream $\dot m_{s,z}$ mix with entrained background air $\dot m_{e,z}$ in a small region around the passenger. With negligible local air storage, balanced outflow, and $x\in\{h,w\}$,
\begin{equation}
 0=\dot m_{s,z}(x_s^{\rm sup}-x_{\ell,s,z})+
 \dot m_{e,z}(x_b-x_{\ell,s,z}),\qquad
 \beta_{s,z}=\frac{\dot m_{s,z}}{\dot m_{s,z}+\dot m_{e,z}}.
\label{eq:delivery-balance}
\end{equation}
A fraction $\chi$ of device electricity may dissipate into the supply air before reaching people. Write $p_s$ for unit power in kW and $P_s=1000p_s$ for watts. With unit dry-air supply flow $\dot m_s$, define
\begin{equation}
 h_s^{\rm sup}=h_{o,s}+\chi P_s/\dot m_s,\qquad
 w_s^{\rm sup}=w_{o,s},\qquad 0\leq\chi\leq1.
\label{eq:supply-heat}
\end{equation}
For a fan, the unheated outlet equals its inlet. Solving Eq.~\eqref{eq:delivery-balance} gives the shared operator
\begin{align}
 w_{\ell,s,z}&=(1-\beta_{s,z})w_b+\beta_{s,z}w_s^{\rm sup},\\
 h_{\ell,s,z}&=(1-\beta_{s,z})h_b+\beta_{s,z}h_s^{\rm sup},\\
 T_{\ell,s,z}&=\frac{h_{\ell,s,z}-2501000w_{\ell,s,z}}{1006+1860w_{\ell,s,z}}.
\label{eq:mixing}
\end{align}
RH is recovered from $T_\ell,w_\ell$, not averaged. All local states must satisfy psychrometric and SET-domain limits. The bus base case uses $b=a$ and $\chi=0$. Its assumed $\beta$ values imply entrainment-to-supply ratios $(1-\beta)/\beta$: $0.25/1.22/9.00$ for B and $0.11/0.43/1.50$ for C. These are an explicit flow-ratio closure, not newly measured airflow. If the strips are supplied in parallel, their supply branches must sum to one device's flow; no full-device flow is counted independently in every strip.

The local quasi-steady approximation requires air residence time to be short relative to weather or operating changes. No measured residence time is available, so the bus calculation is a conditional steady-state comparison. Its fixed outdoor reservoir also omits accumulated crowd heat and humidity. The station extension releases that approximation while retaining Eqs.~\eqref{eq:outlet}--\eqref{eq:mixing} and the response below.

\subsection{Converting local conditions to a common human response}

We compute standard effective temperature (SET), denoted $\Theta$, from
\begin{equation}
 \Theta_{s,z}=F_{\mathrm{SET}}(T_{\ell,s,z},T_{r,z},v_{s,z},\phi_{\ell,s,z};\mathrm{met},\mathrm{clo}).
\label{eq:set}
\end{equation}
SET compares physiological heat exchange with a reference environment, incorporating temperature, humidity, air movement and radiation \cite{thermalcomfort,fanguide}. We use \texttt{pythermalcomfort} with applicability checks \cite{pythermal}. SET is a quasi-steady engineering proxy here; it does not resolve each passenger's prior walking or transient physiology.

\subsection{Measuring cooling efficiency}

For the warm-weather scenarios, define the hot-discomfort proxy
\begin{equation}
 D(\Theta)=[\Theta-\Theta_c]_+,\qquad [x]_+=\max(x,0).
\end{equation}
The threshold is an assumed reference defining useful relief, not a scheme-admission constraint or an ASHRAE safety limit. Let $0$ denote the same shelter without active cooling, retaining the same shade, weather, radiation, and background air movement. The signed local improvement is
\begin{equation}
 b_{s,z,k}=D(\Theta_{0,z,k})-D(\Theta_{s,z,k}).
\end{equation}
Negative values are retained: an adverse condition must reduce useful service. Over time steps of $\Delta t_k$ hours,
\begin{align}
 U_s&=\sum_k\sum_z n_{z,k}b_{s,z,k}\Delta t_k, &
 E_s&=\sum_k p_{s,k}\Delta t_k,\\
 \eta_s&=U_s/E_s,\qquad E_s>0.
\label{eq:bus-efficiency}
\end{align}
The unit of $\eta$ is person$\cdot{}^{\circ}\mathrm C\cdot$h/kWh, with temperature measured on the SET scale. This is a service-efficiency index, not a thermodynamic percentage or COP. Empty zones earn no credit, although operating electricity is still counted. The zero-energy baseline has no efficiency ratio.

For service accounting, define remaining discomfort,
\begin{equation}
 R_s=\sum_k\sum_z n_{z,k}D(\Theta_{s,z,k})\Delta t_k,
 \qquad \overline D_s=\frac{R_s}{\sum_k\sum_z n_{z,k}\Delta t_k}.
\label{eq:remaining}
\end{equation}
The identity $U_s+R_s=R_0$ checks that relief is counted against the same baseline. Here $R_s$ is a diagnostic of the numerator, not an additional ranking objective: Requirements 2 and 3 compare $\eta_s$. Its later use as an adequacy objective or constraint is introduced in Section~\ref{sec:effectiveness}. A comfortable baseline earns no additional hot-discomfort relief.

\noindent\textbf{Geometry and demand.} Use either headcounts $n_z$ or gross-area densities $\rho_z$ with $n_z=\rho_z A_z$. Area is counted once, not multiplied into the numerator again. The uniform and eight-person-strip cases imply $1.85$ and $3.69$ persons per gross square metre; these are concentration stress scenarios, not demonstrated usable capacities. Bench and circulation dimensions are needed for a deployable layout. Fixed headcounts and fields give fixed service; fixed densities and fields give service proportional to area. These are accounting properties, not a claim that airflow is independent of geometry.

\noindent\textbf{A ranking result independent of the assumed delivery maps.} Under the interpreted constant powers and a common operating period,
\begin{equation}
\begin{aligned}
 \eta_B>\eta_A&\iff U_B>4.8U_A,\\
 \eta_C>\eta_A&\iff U_C>10U_A,\\
 \eta_C>\eta_B&\iff U_C>\frac{25}{12}U_B.
\end{aligned}
\label{eq:break-even}
\end{equation}
Thus nominal temperature reduction alone cannot justify a larger device. Its accumulated passenger relief must compensate for its larger electricity input. These inequalities follow from the given power ratios; their validity does not depend on the assumed zone profiles.

\subsection{Numerical results: mechanisms and efficiency ranking}

\begin{figure}[!ht]
\centering
\includegraphics[width=0.90\linewidth]{figures/01_local_conditions.pdf}
\caption{Simulated local air temperature, RH, and SET under the stated assumptions. Fan A leaves air temperature unchanged while changing SET. B and C lower air temperature and increase RH. These are scenario outputs, not field measurements.}
\label{fig:bus-local}
\end{figure}

Spatial delivery matters even within one small shelter. In the dry case, C gives SET $26.3\,^{\circ}\mathrm C$ near the outlet but $29.4\,^{\circ}\mathrm C$ at the edge. At RH $70\%$, these become $33.1$ and $35.2\,^{\circ}\mathrm C$. A single mean air-temperature reduction would conceal these differences. A, meanwhile, improves SET without changing the $34\,^{\circ}\mathrm C$ dry-bulb temperature.

\begin{figure}[!ht]
\centering
\includegraphics[width=0.90\linewidth]{figures/24_bus_efficiency.pdf}
\caption{Uniform occupancy: cooling efficiency $\eta=U/E$ for the two assumed humidity cases. SET-based relief defines the numerator; the ranking criterion is efficiency alone.}
\label{fig:bus-efficiency}
\end{figure}

For uniform occupancy, dry-case efficiencies are $84.6$, $24.5$ and $21.1$ for A, B and C. At RH $70\%$, they are $90.6$, $21.1$ and $15.2$. Thus A is the efficiency leader under both specified climates. Requirements 2 and 3 retain this ratio-based comparison; Section~\ref{sec:effectiveness} considers what these ratios do not establish about overall service.

\subsection{Climate and passenger distribution}

\begin{figure}[!ht]
\centering
\includegraphics[width=0.90\linewidth]{figures/03_humidity_sweep.pdf}
\caption{RH sweep with temperature, radiation, occupancy, and delivery profiles fixed. Outlet cooling is bounded by both the specified cap and the wet-bulb depression.}
\label{fig:bus-humidity}
\end{figure}

B/C's assumed outlet drops are $5/8\,^{\circ}\mathrm C$ at RH $30\%$ but approximately $3.4/3.4\,^{\circ}\mathrm C$ at RH $70\%$. Equal high-humidity drops follow from the common assumed effectiveness once neither cap binds; delivery profiles still differ. Nonlinear SET and a changing baseline mean service efficiency need not vary monotonically with RH.

In dry conditions, A's efficiency falls from $118.3$ with near-heavy occupancy to $48.3$ with edge-heavy occupancy. B exceeds C under uniform occupancy ($24.5$ versus $21.1$), but C exceeds B under edge-heavy occupancy ($16.6$ versus $15.5$). A remains the efficiency leader in all six nominal climate--occupancy cases. Full occupancy and endpoint-sensitivity figures remain separately available with the computational material.

\subsection{When the efficiency ranking changes}

A poorly aimed or obstructed fan can consume electricity without delivering its nominal air movement to passengers. We test this specific failure mode using
\begin{equation}
 v_{A,z}(q)=v_0+q(v_{A,z}^{\mathrm{nom}}-v_0),\qquad 0\leq q\leq1,
\end{equation}
while holding $p_A=0.25$ kW. Here $q$ describes delivery loss at occupied positions, \emph{not} a fan-speed control setting; changing motor speed would require a different power model. The B/C profiles stay fixed in this particular stress test.

\begin{figure}[!ht]
\centering
\includegraphics[width=0.90\linewidth]{figures/05_delivery_decision.pdf}
\caption{Conditional efficiency winner on the complete RH--fan-delivery grid: RH $20$--$80\%$ in one-percentage-point steps and $q=0$--$1$ in steps of $0.02$. The boundary is grid-resolved. This tests fan delivery loss; it is not a universal climate-selection map.}
\label{fig:bus-decision}
\end{figure}

The first sampled $q$ at which A wins is $0.18$ at RH $30\%$ and $0.16$ at RH $70\%$; the switch lies between that sample and the preceding $0.02$ step. B wins below these levels; C wins no cell of this specific grid. These are conditional delivery thresholds, not universal climate recommendations. Measure occupied-location $(T,\mathrm{RH},v,T_r)$, electricity and occupancy to replace the assumed profiles before using the comparison for an installation.

\subsection{Sensitivity, verification, and scope}



Across the tested individual endpoint changes, the largest B/A and C/A efficiency ratios are $0.432$ and $0.372$; no endpoint reverses the nominal advantage of A. This does not prove robustness over every intermediate value or joint parameter combination. Added-speed factors scale $v-v_0$, and cooled-air fractions are capped at one. The rating-interpretation test removes the $5/8\,^{\circ}\mathrm C$ outlet caps while retaining the wet-bulb constraint; it addresses uncertainty in where the quoted temperature effects were measured. The delivery-loss test in Figure~\ref{fig:bus-decision} already demonstrates that stronger departures can change the winner.

\noindent\textbf{Relaxing the bus heat approximation.} A separate $\chi=0,0.5,1$ test assumes unit supply flows $0.6/0.6/1.2$ kg dry air/s and, for A's previously immaterial scalar mixing, $\beta_A=(0.8,0.45,0.1)$. These additional assumptions are not inferred from rated power. From $\chi=0$ to 1, dry-case A/B/C efficiencies change from $84.6/24.5/21.1$ to $81.0/20.8/18.0$; humid-case values change from $90.6/21.1/15.2$ to $86.8/17.4/12.1$. A remains the most efficient in these tested cases. This quantifies the efficiency consequence of the introductory zero-local-heat approximation; the rail comparison includes this heat and tests its allocation.

\noindent\textbf{Verification.} The implementation checks psychrometric round trips, dry-air moisture/enthalpy mixing balances, wet-bulb bounds, and saturation. Service-accounting checks verify $U_s+R_s=R_0$, zero benefit at zero occupancy, zero benefit when all delivered effects vanish, and unchanged rankings when every zone's population doubles. A signed-adverse-effect check confirms that worsening discomfort is not clipped away. No credit is given for additional cooling of an already comfortable baseline. Geometry checks verify unit conversion and area conservation, equivalence of count and density inputs, fixed-count invariance, and fixed-density scaling under unchanged local fields.

\Needspace{4\baselineskip}
\noindent\textbf{Limitations and transfer.} Local profiles and representative passengers remain assumptions; the numerical ranking is conditional on them. The conservation-based mixing closure uses assumed flow ratios; fan curves, crosswind and opening geometry have not been used to identify them. Radiation is fixed for a shaded shelter, so direct sun and changing shadow patterns need additional modeling. Short waiting times and varied physiological histories are not resolved by quasi-steady SET. Without measured flow, water consumption cannot be calculated reliably from the moisture increase alone; water and operating cost belong in the later effectiveness analysis. Section~\ref{sec:rail} retains the same device, local mixing, SET and service modules, and adds a dynamic background, surface storage and variable equipment deployment.

\noindent\textbf{Reproducibility.} The project retains PsychroLib 2.5.0 and \texttt{pythermalcomfort} 4.8.1 inputs, source code, CSV results, software versions and verification records. Each figure is a separately replaceable PDF. These are conditional simulations, not experimental validation.

\clearpage
\section{Canberra MRT Station: A Semi-open Waiting-area Case}
\label{sec:rail}

\subsection{Site evidence and the modeled waiting area}

We select Canberra MRT station in Singapore for Requirement 2(b). RSP describes an elevated, two-storey station with a site area of $9,128\,\mathrm{m^2}$ \cite{canberra-rsp}. LTA identifies platform planting, daylight-admitting glazing and rain-protective louvers that allow natural ventilation \cite{canberra-lta}. The efficiency question is how device type, quantity and location change useful passenger relief per kWh. Coverage enters service weighting; overall service requirements and final installation choice are addressed in Section~\ref{sec:effectiveness}. We evaluate candidate installations of A, B and C in this shaded, humid, semi-open setting; their deployment is hypothetical.

\begin{figure}[!ht]
\centering
\includegraphics[width=0.84\linewidth]{figures/13_canberra_station.pdf}
\caption{User-supplied Canberra station schematic, retained as a separate PDF. Approximate structural dimensions are contextual references; the drawing is not an as-built survey or a map of the calculation zones. Sources identified in the drawing are listed in \cite{canberra-dimensions,canberra-rsp}.}
\label{fig:canberra}
\end{figure}

\noindent\textbf{Modeled occupied area.} The schematic labels approximately 220 m length, roof span over 28 m and roughly 17 m columns. Roof span and column height are not occupied width or platform clearance. We assume aggregate waiting width $b=6$ m, usable fraction $\alpha=0.8$ and effective background-air height 4 m:
\begin{equation}
 A_{\mathrm{occ}}=\alpha Lb=1056\,\mathrm{m^2}.
\label{eq:canberra-area}
\end{equation}
Three equal aggregate zones, access/core/edge, each have $352\,\mathrm{m^2}$ and $1408\,\mathrm{m^3}$. They represent contrasting exposure conditions, not surveyed compartments. Neither occupied area nor effective height is inferred directly from site area or column dimensions.

\subsection{What changes from the bus model, and why}

The service objective, device response, local mixing operator, SET and energy denominator remain unchanged. What changes is the background state supplying each device. The bus uses a prescribed outdoor reservoir; the station solves Eq.~\eqref{eq:common-balance} with time-varying people, exchange and surface feedback. Local delivered air is then computed using the same Eq.~\eqref{eq:mixing}. This makes the extension traceable at both the physical and service levels.

\begin{table}[!ht]
\centering\small
\caption{Retained modules, released approximations, and their consequences.}
\label{tab:rail-changes}
\begin{tabularx}{\linewidth}{@{}>{\raggedright\arraybackslash}p{0.24\linewidth}XX@{}}
\toprule
Module & Bus simplification & Rail extension and purpose\\
\midrule
Device and local delivery & Fixed outdoor inlet; steady entrainment; local electric heat neglected. & Evolving background inlet; same mixing balance with supply heat, to distinguish local relief from background response.\\
Background heat/moisture & Outdoor reservoir is prescribed. & Solve storage and interzone exchange to capture humidity feedback and time lag.\\
Surfaces & Fixed shaded radiant environment. & Solve effective surface storage and radiative influence.\\
Response and service & SET, signed relief, people-time weighting and all-device energy. & Retain all four; vary counts and locations and assess coverage.\\
\bottomrule
\end{tabularx}
\end{table}

The assumed base exchange of $4$--$10$ h$^{-1}$ gives an air replacement time $(3600/ACH)\,\mathrm{s}=6$--$15$ min, comparable to the assumed 15-min passenger cycle. This supports resolving background dynamics in this scenario; station size alone would not justify it. The local mixing region is treated as having negligible storage only if $\tau_\ell=M_\ell/(\dot m_k+\dot m_e)$ is small relative to the 120-s arrival disturbance. $M_\ell$ is the air held in the local jet core, not footprint area multiplied by the platform height. Its unmeasured residence time is a limitation of this closure.

\subsection{Coupled heat, moisture, and surface balances}

Following the control-volume structure of building heat and moisture balances \cite{epzone}, let each zone store moist-air enthalpy $h_z$ (J/kg dry air), humidity ratio $w_z$ (kg/kg dry air), and an effective roof/wall inner-surface temperature $T_{m,z}$. The background air is well mixed within each aggregate zone; the embedded local delivery regions are evaluated separately with negligible air storage. Background dry-air mass is approximated as constant, $M_z=\rho_{da,o}V_z$, with outdoor reference density. Using the occupied-zone volume for background storage neglects the smaller local jet hold-up. This is a reduced model, not an EnergyPlus or CFD simulation.

Define the exchange operator, for either $x=h$ or $w$, as
\begin{equation}
 \mathcal L_z(x)=\dot m_{0z}(x_o-x_z)+\sum_{j\ne z}\dot m_{zj}(x_j-x_z),
 \qquad \dot m_{0z}=M_z\lambda_z/3600.
\label{eq:rail-exchange}
\end{equation}
Here $\lambda_z$ is the assumed outdoor air-change rate in h$^{-1}$. Each exchange represents equal entering and leaving dry-air flow; adjacent-zone rates satisfy $\dot m_{zj}=\dot m_{jz}$. Wind, buoyancy and opening geometry can affect outdoor exchange \cite{ventilation}; a prescribed train-cycle pulse tests that uncertainty without claiming a station-specific piston-flow law.

With $x_{zk}$ devices of type $k$ in zone $z$, unit circulation $\dot m_k$ (kg dry air/s), and time in seconds, our balances are
\begin{align}
 M_z\dot h_z={}&\mathcal L_z(h)+H_{ma,z}(T_{m,z}-T_z)
 +n_zq_c+n_zg h_v(T_z)+P_z
 +\sum_k x_{zk}\dot m_k(h_{e,zk}-h_z),\label{eq:rail-enthalpy}\\
 M_z\dot w_z={}&\mathcal L_z(w)+n_zg+
 \sum_k x_{zk}\dot m_k(w_{e,zk}-w_z),\label{eq:rail-moisture}\\
 C_{m,z}\dot T_{m,z}={}&H_{mo,z}(T_o-T_{m,z})+Q_{\mathrm{sol},z}
 +n_zq_r-H_{ma,z}(T_{m,z}-T_z).\label{eq:rail-surface}
\end{align}
The per-person inputs are convective heat $q_c$, radiative heat $q_r$, and water-vapor generation $g$. Vapor enters with $h_v(T)=2501000+1860T$ J/kg; no additional latent-heat term is added. With $P_k=1000p_k$ W per unit, total device power is $P_z=\sum_k x_{zk}P_k$ (W). This power ultimately dissipates within the modeled zone. Its division between supply air and background affects local exposure but not total input. Pump power is already included in the device total. Conductances $H$ have units W/K, surface capacity $C_m$ has units J/K, and $Q_{\mathrm{sol}}$ is an effective absorbed envelope load. The opposite signs of the air--surface term conserve internal energy.

Air temperature is recovered from $T_z=(h_z-2501000w_z)/(1006+1860w_z)$. A shaded longwave approximation for the surface contribution to comfort \cite{thermalcomfort} gives
\begin{equation}
 T_{r,z}=\left[F_z(T_{m,z}+273.15)^4+(1-F_z)(T_{r,o}+273.15)^4\right]^{1/4}-273.15.
\label{eq:rail-mrt}
\end{equation}
$F_z$ is an assumed effective view fraction; the example sets external radiant temperature $T_{r,o}=T_o$, while keeping it a separate input. Direct sun on passengers and detailed radiative asymmetry are excluded.

\noindent\textbf{The retained device module.} Evaluate Eq.~\eqref{eq:outlet} using the current background inlet. For A, the unheated circulation outlet equals the background, so its circulation has no net bulk enthalpy or water source. For B/C, $h_{e,zk}\simeq h_z$ while $w_{e,zk}>w_z$ when evaporation occurs. Increasing background moisture reduces subsequent wet-bulb depression. Here the subscript $e$ denotes the unheated outlet, the state labeled $o,s$ in Eqs.~\eqref{eq:outlet}--\eqref{eq:enthalpy}; it does not denote outdoor air. Device electricity remains a separate positive heat source in Eq.~\eqref{eq:rail-enthalpy}.

\subsection{Local delivery without double-counting heat or water}

For each covered subgroup, apply Eqs.~\eqref{eq:supply-heat}--\eqref{eq:mixing} with background $b=z$. We distinguish the treated-stream fraction $\beta_{zk}$ from the passenger coverage fraction $c_{zk}$ defined below. The supply flow draws from the background, entrains additional background air, and returns the mixed air to that same zone. For finite entrainment ($\beta>0$),
\begin{align}
 (\dot m_k+\dot m_e)(h_{\ell,zk}-h_z)+(1-\chi)P_k
 &=\dot m_k(h_{e,zk}-h_z)+P_k,\\
 (\dot m_k+\dot m_e)(w_{\ell,zk}-w_z)
 &=\dot m_k(w_{e,zk}-w_z).
\label{eq:return-closure}
\end{align}
Thus the return flow gives exactly the existing background source. The local state improves the exposure calculation without adding evaporation or electrical heat a second time. People outside direct footprints experience the background state; those inside experience the mixed supply state. Both experience the zone radiant environment. Prescribed air speeds remain a separate approximation; no momentum or detailed jet trajectory is solved.

We set $\beta=0.5$ for each type as an explicit equal-entrainment scenario and use unit circulation flows $0.6/0.6/1.2$ kg dry air/s for A/B/C. These are assumptions, not manufacturer data. The nominal $\chi=1$ puts all electrical dissipation into the air before it reaches passengers, a conservative supply-heating case. The sensitivity analysis varies $\beta$ and $\chi$, and separately varies A's newly assumed circulation flow. The case $\beta=0$ is an infinite-dilution scalar limit that exactly recovers the previous shared-background exposure; it retains the prescribed local speeds and is a diagnostic limit, not a measured finite-flow state.

\noindent\textbf{Recovery of the bus model.} Hold the background at outdoor $h_a,w_a$, prescribe the bus radiant environment, speeds and occupancies, neglect local electrical heating ($\chi=0$), and use the bus $\beta$ values. The same device and delivery functions then give the bus temperatures, humidities, SET and service exactly. Fixing the background corresponds to an outdoor-reservoir approximation to Eq.~\eqref{eq:common-balance}; merely setting time derivatives to zero would not impose it. This distinction makes the reduction physically explicit.

\subsection{Device counts, allocation, and service weighting}

Here $s$ indexes an installation and $k$ a device type. Let $x_{zk}\in\mathbb Z_{\ge0}$ be the number of devices of type $k\in\{A,B,C\}$ in zone $z$. A $3\times3$ count matrix specifies a candidate installation, allowing quantities and device types to vary by zone. Each candidate is evaluated through the resulting shared thermal environment: all operating units contribute heat and electricity, and evaporative units also contribute moisture.

We use assumed, non-overlapping service footprints $a_A=20$, $a_B=30$ and $a_C=40\,\mathrm{m^2}$ per unit. These are comparison inputs, not manufacturer coverage ratings. For uniform passenger density within a zone, define
\begin{equation}
 c_{zk}=\frac{x_{zk}a_k}{A_z},\qquad
 c_{z0}=1-\sum_k c_{zk}\ge0,\qquad \sum_kx_{zk}a_k\le A_z.
\label{eq:coverage}
\end{equation}
The uncovered fraction experiences local speed $v_0$; each covered subgroup experiences its device-specific speed $v_{zk}$. We average \emph{discomfort}, not velocities or device efficiencies:
\begin{equation}
 \overline D_{s,z}(t)=c_{z0}D[\Theta(T_z,T_{r,z},v_0,\phi_z)]
 +\sum_k c_{zk}D[\Theta(T_{\ell,zk},T_{r,z},v_{zk},\phi_{\ell,zk})].
\label{eq:mixed-discomfort}
\end{equation}
The subgroups share radiant conditions but can differ in air temperature, humidity and speed. An uncovered passenger can still experience indirect background cooling or humidification. Coverage denotes direct delivery, not a boundary of all thermal influence. Separate, non-overlapping footprints approximate distributed equipment; co-located jets, overlap and series operation require additional airflow modeling.

\subsection{Canberra scenario and comparable deployment tests}

We assume $33\,^{\circ}\mathrm C$ and $70\%$ outdoor RH for a hot, humid Canberra scenario. Singapore's weather service describes typical daytime maxima of $31$--$33\,^{\circ}\mathrm C$ and high humidity, with rain-free afternoon RH around $60\%$ \cite{singapore-climate}. Our pair is a scenario, not a simultaneous observation at the station. SET retains $1.2$ met, $0.5$ clo, and the common $26\,^{\circ}\mathrm C$ reference.

\begin{table}[!ht]
\centering\small
\caption{Station-scale scenario: operating assumptions, with powers and outlet caps interpreted from the problem.}
\label{tab:rail-inputs}
\begin{tabularx}{\linewidth}{@{}>{\raggedright\arraybackslash}p{0.27\linewidth}>{\raggedright\arraybackslash}X@{}}
\toprule
Input & Assumed value and interpretation\\
\midrule
People and train cycle & An assumed 15-min stress-test cycle, not the actual timetable. Zone counts rise from $(26.4,44,17.6)$ to $(88,198,66)$, then reset; fractional values represent mean counts.\\
Outdoor and zone exchange & Base ACH $(8,4,10)$ h$^{-1}$, plus $(8,4,12)$ in the last 2 min; adjacent-zone exchange $0.88$ kg dry air/s.\\
Surface response, per zone & $C_m=17.6$ MJ/K; $H_{ma}=880$, $H_{mo}=1760$ W/K; absorbed load $3520$ W; view fractions $(0.55,0.85,0.55)$.\\
Human sources & Convection $45$ W, radiation $30$ W, vapor $45/(2.45\times10^6)$ kg/s per person.\\
Device powers and flows & $p_{A/B/C}=0.25/1.2/2.5$ kW; circulation $0.60/0.60/1.20$ kg dry air/s per unit. B/C use $\varepsilon=0.70$ and $5/8\,^{\circ}\mathrm C$ outlet-drop caps. Local $\beta=0.5$, $\chi=1$.\\
Speeds at covered locations & A $(0.60,1.00,0.50)$, B $(0.50,0.80,0.45)$, C $(0.75,1.00,0.60)$ m/s; uncovered locations $0.20$ m/s.\\
\bottomrule
\end{tabularx}
\end{table}

The same occupied area, thermal properties, load densities and passenger schedule are used for every candidate. This holds demand constant while equipment quantities and allocation change. The architectural sources establish the venue context; the assumed operating conditions and coverage parameters define the numerical comparison.

We first compare six benchmark configurations at the same direct-service footprint: $240\,\mathrm{m^2}$ per zone, or $68.2\%$ of the assumed occupied area. Per-zone counts $(x_A,x_B,x_C)$ are $(12,0,0)$, $(0,8,0)$, $(0,0,6)$, $(6,4,0)$, $(6,0,3)$ and $(3,2,3)$. Whole-station totals are respectively 36, 24, 18, 30, 27 and 24 units. Device-count sweeps test how efficiency changes with deployment scale; a fixed-inventory redistribution test isolates allocation. Equal-coverage benchmarks control delivery extent without imposing a minimum-service requirement.

A prescribed two-hour uncooled warm-up gives all candidates the same initial air and surface states. Each then runs for two hours alongside a no-device trajectory with identical weather, people and exchange. This specifies common prior history without claiming periodic steady state; no free initial stored cooling is credited. Electricity covers all cooling devices, not the whole station bill; benefits after the horizon are excluded.

\subsection{Comparing cooling efficiency}

Retain signed relief against the matched baseline, now using the coverage-weighted discomfort:
\begin{align}
 U_s&=\frac1{3600}\int_0^{7200}\sum_z n_z(t)[D_{0,z}(t)-\overline D_{s,z}(t)]\,dt,\\
 E_s&=\frac1{3600}\int_0^{7200}\sum_{z,k}x_{zk}p_k\,dt,
 \qquad \eta_s=U_s/E_s.
\label{eq:rail-efficiency}
\end{align}
With $t$ in seconds and $p_k$ in kW, the factors $1/3600$ convert to hours. The unit is person$\cdot{}^{\circ}\mathrm C\cdot$h/kWh. Compare the candidates by $\eta_s$ under the common scenario; no residual-discomfort, coverage or budget acceptance thresholds are imposed here. Coverage remains part of the service numerator through the subgroup weights. Mean remaining discomfort $\overline R_s=\int\sum_z n_z\overline D_{s,z}\,dt/\int\sum_z n_z\,dt$ is retained in the output as a diagnostic for the later effectiveness extension.

% UPDATED_RESULTS_START
The common two-hour comparison contains 440 person-hours. We distinguish the full set of 24 sampled plans from the six controlled benchmarks with equal $68.2\%$ direct-service coverage. In the full set, A3 has the highest $\eta$, $37.58$ person$\cdot{}^{\circ}$C$\cdot$h/kWh. Within the equal-coverage subset, A36 is highest at $35.49$. A36 is not the efficiency maximum across all counts. These are conditional maxima over finite candidate sets, not global optima.

\begin{table}[!ht]
\centering\small
\caption{Efficiency accounting under the nominal assumptions ($\beta=0.5$, $\chi=1$). $U$ is useful relief in person$\cdot{}^{\circ}$C$\cdot$h; $E$ is kWh. Coverage describes the service weights, not a separate objective.}
\label{tab:deployment-results}
\begin{tabular}{@{}lrrrr@{}}
\toprule
Plan & Person-time coverage & $U$ & $E$ & $\eta=U/E$\\
\midrule
A3 & 5.68\% & 56.38 & 1.5 & 37.58\\
A36 & 68.18\% & 638.89 & 18.0 & 35.49\\
B24 & 68.18\% & 446.28 & 57.6 & 7.75\\
C18 & 68.18\% & 453.00 & 90.0 & 5.03\\
AB30 & 68.18\% & 581.61 & 37.8 & 15.39\\
AC27 & 68.18\% & 593.68 & 54.0 & 10.99\\
ABC24 & 68.18\% & 524.76 & 63.9 & 8.21\\
AB30-R & 76.19\% & 695.07 & 37.8 & 18.39\\
\bottomrule
\end{tabular}
\end{table}

\noindent\textbf{Scale and allocation.} Increasing the number of fans increases served person-time, while their heat also changes the background; the added useful relief need not be proportional to electricity. The resulting count sweep is evaluated through $U/E$. A3's $5.68\%$ coverage is recorded as the scope of its service, without imposing a coverage threshold in this efficiency comparison.

AB30-R redistributes the same 18 A and 12 B units from equal counts to $(3,4,0)$, $(12,3,0)$ and $(3,5,0)$. Electricity remains 37.8 kWh and area coverage $68.2\%$, while person-time coverage increases from $68.2\%$ to $76.2\%$. Useful relief rises from $581.61$ to $695.07$ person$\cdot{}^{\circ}$C$\cdot$h, increasing efficiency from $15.39$ to $18.39$. The gain results from delivering the same inventory's cooling closer to passenger demand.

\noindent\textbf{Efficiency sensitivity.} For the 15 tested combinations $\beta=0,0.25,0.5,0.75,1$ and $\chi=0,0.5,1$, A36 remains the efficiency leader among the six equal-coverage benchmarks. Changing A's circulation from $0.3$ to $0.6$ to $1.2$ kg dry air/s gives A36 efficiencies $34.06/35.49/36.22$; it retains its equal-coverage efficiency lead. The separate footprint-sensitivity figure reports efficiency changes at fixed equipment counts. These are defined parameter tests, not robustness over every possible climate or delivery pattern. Overall adequacy and tradeoffs are considered separately in Section~\ref{sec:effectiveness}.

% UPDATED_RESULTS_END

\begin{figure}[!ht]
\centering
\includegraphics[width=0.90\linewidth]{figures/22_station_efficiency.pdf}
\caption{Efficiency comparisons from the same frozen simulations. A3 leads all 24 sampled plans; A36 leads the six equal-coverage benchmarks. Device-count curves show how the ratio changes with scale.}
\label{fig:deployment-efficiency}
\end{figure}





\begin{figure}[!ht]
\centering
\includegraphics[width=0.96\linewidth]{figures/23_delivery_efficiency.pdf}
\caption{Efficiency sensitivity to local stream mixing and supply-heat allocation. All cases use the same background trajectories; $\beta=0$ recovers shared-background exposure. No additional objective is optimized here.}
\label{fig:local-sensitivity}
\end{figure}

\subsection{Verification, evidence, and transfer}

The implementation checks both the full background budgets and the local return identities in Eq.~\eqref{eq:return-closure}. Limiting tests recover the old shared-background station calculation at $\beta=0$ and the bus calculation under its fixed-background assumptions. Refining the time step tests integration error; varying $\beta,\chi$ changes local exposure without changing the already balanced background trajectories. The complete check results and trajectories accompany the paper.



The experimental study that informed this structure reports an air-temperature MAE of $0.4\,^{\circ}\mathrm C$ and an RH MAE of about five percentage points for its mist-on validation period \cite{cruz2026}. It separates calibration and validation periods and compares operation against a matched no-mist simulation. We adopt those comparison principles, but its dry-season, unoccupied cafeteria is not a calibration dataset for Canberra. Its reported average $1.9\,^{\circ}\mathrm C$ cooling and mist effectiveness coefficients are not copied into our scenario. It also did not evaluate occupant comfort, so its cooling effectiveness cannot replace our SET-based service measure.

For field validation, first measure the no-device background and estimate exchange and surface response; then measure device outlet and occupied-location $T$, RH and speed to estimate delivery and heat allocation. Use separate periods to test predictions against matched baseline conditions and metered electricity. Compare occupied positions as well as background averages: the study's own well-mixed-zone limitation motivates our local-delivery refinement. Device footprints, passenger distribution, residence times and flows remain uncertain, so present results are conditional comparisons of sampled installations.

The model excludes direct sun, detailed stratification, surface moisture sorption and condensation; it stops on invalid humidity or SET inputs. SET remains quasi-steady even when background conditions evolve. Source code, scenario inputs, numerical checks, CSV outputs and separately replaceable PDF figures are retained in the project. The same chain can be applied to later venues by changing evidenced boundaries and inputs while retaining device, delivery, response and service modules.


\clearpage
\section{Overall Cooling Effectiveness (Requirement 4)}
\label{sec:effectiveness}

Requirements 2 and 3 compare useful cooling per kWh, $\eta=U/E$. Requirement 4 asks a different decision question: which installation balances cooling performance, visitor experience, energy and cost? The physical model and its outputs are retained; the decision criteria are expanded. This section organizes that extension. Cross-venue applications for Requirement 3 and empirical cost/experience inputs are still pending in this working draft.

\subsection{Four objectives with distinct roles}

For a common operating and accounting horizon $H$, compare the vector
\begin{equation}
 \min_s\big(\overline R_s,\ X_s,\ E_{s,H},\ C_{s,H}\big).
\label{eq:effectiveness-vector}
\end{equation}
Here $\overline R_s$ is mean remaining SET excess, so lower values indicate more complete thermal relief. $X_s$ represents non-thermal visitor-experience loss, such as the fraction reporting unacceptable equipment noise or obstructed access; it requires observations or surveys. SET-based thermal discomfort is not added to $X_s$ again. $E_{s,H}$ is energy use and $C_{s,H}$ is monetary cost over the same horizon. Because $U_s+R_s=R_0$ under a matched baseline, useful relief and residual discomfort are not treated as independent objectives, and $\eta$ is not added as a fifth objective.

A transparent cost model is
\begin{equation}
 C_{s,H}=\sum_k N_{s,k}\big(c_{k,H}^{\rm own}+c_{k,H}^{\rm maint}\big)
 +\pi_e E_{s,H}+\pi_w V_{s,H},
\label{eq:total-cost}
\end{equation}
where $N_{s,k}=\sum_zx_{zk}$, ownership and maintenance costs are allocated to $H$, $\pi_e,\pi_w$ are energy and water prices, and $V_{s,H}$ is make-up water including losses. All costs use one currency; $\pi_e$ is currency/kWh, $\pi_w$ is currency/m$^3$ and $V$ is m$^3$. Prices, ownership costs, water losses and experience data have not been supplied. We therefore retain these as inputs rather than inventing a calibrated cost or experience ranking. Energy represents a resource objective, while cost represents expenditure; the correlated quantities are not simply added with arbitrary equal weights.

\subsection{Service requirements and a conditional decision rule}

Let $q_s$ be person-time coverage, already used to account for exposure in Requirement 2. In this extension an operator may additionally require $q_s\ge q_{\min}$ and $\overline R_s\le R_{\max}$. For an electricity-priority policy, the earlier selection rule now belongs here:
\begin{equation}
 s^*=\arg\min_s E_{s,H}\quad\text{subject to}\quad
 \overline R_s\le R_{\max},\quad q_s\ge q_{\min},\quad
 X_s\le X_{\max},\quad C_{s,H}\le B_H.
\label{eq:effectiveness-selection}
\end{equation}
The bounds express stated service, experience and budget requirements, not universal SET safety limits. A cost-priority operator could instead minimize $C_{s,H}$ with an energy cap. If no priorities or bounds are specified, report nondominated choices for Eq.~\eqref{eq:effectiveness-vector}; no unique overall winner follows from efficiency alone.

\subsection{What the existing simulations demonstrate}

Across the same full set of 24 station candidates, A3 maximizes efficiency, while A36 delivers much wider direct service (Table~\ref{tab:efficiency-effectiveness}). These observations motivate the extension; they do not establish a complete effectiveness ranking.

\begin{table}[!ht]
\centering\small
\caption{Efficiency and service adequacy are different questions. Nominal two-hour simulations, with the same weather and passenger demand.}
\label{tab:efficiency-effectiveness}
\begin{tabular}{@{}lrrrr@{}}
\toprule
Plan & $\eta$ & Person-time coverage & $\overline R_s$ ($^{\circ}$C) & $E$ (kWh)\\
\midrule
A3 & 37.58 & 5.68\% & 10.44 & 1.5\\
A36 & 35.49 & 68.18\% & 9.12 & 18.0\\
\bottomrule
\end{tabular}
\end{table}

For example, $q_{\min}>q_{\mathrm{A3}}$ excludes A3, where $q_{\mathrm{A3}}\approx5.68\%$. A36 passes the coverage test only when $q_{\min}\le q_{\mathrm{A36}}$, where $q_{\mathrm{A36}}\approx68.18\%$, and also requires an acceptable residual, experience and cost. This is not a two-plan optimization: all 24 candidates must be screened under the same requirements.

\begin{figure}[!ht]
\centering
\includegraphics[width=0.86\linewidth]{figures/14_deployment_tradeoff.pdf}
\caption{Requirement 4: the quantified energy--remaining-discomfort projection. The dashed frontier concerns these two outputs of the 24 sampled plans. Experience and cost must be incorporated before claiming a four-objective frontier or an overall recommendation.}
\label{fig:effectiveness-tradeoff}
\end{figure}

The nominal energy--discomfort frontier contains A3, A9, A18, A27, A36, A45 and A51. Delivery uncertainty also matters to the performance objective: among the six equal-coverage benchmarks, AC27 leaves the least discomfort when $\chi=0$ and $\beta\ge0.5$, or when $(\beta,\chi)=(1,0.5)$, although A36 retains the highest efficiency in all 15 tested combinations. At $(0.5,0)$ their residuals are $9.02$ and $9.06\,^{\circ}\mathrm C$, while electricity is 54 and 18 kWh. The full performance/efficiency sensitivity remains available as a separate PDF in the computational material. These tradeoffs belong to overall effectiveness; the visitor-experience and cost inputs must be completed before making the final operating recommendation.


\clearpage
\phantomsection
\addcontentsline{toc}{section}{References}
\begin{thebibliography}{99}
\fontsize{9.5}{11}\selectfont
\urlstyle{same}
\setlength{\itemsep}{0pt}
\setlength{\parskip}{0pt}
\bibitem{ashrae}
ASHRAE. \emph{2026 ASHRAE Handbook---Refrigeration},
Chapter 15, \emph{Retail Food Store Refrigeration and Equipment}.
Product Temperatures and Display Refrigerator Cooling Load sections.
\url{https://handbook.ashrae.org/Handbooks/R26/IP/r26_ch15/r26_ch15_ip.aspx}.
Accessed November 7, 2026.

\bibitem{etl}
UK Department for Energy Security and Net Zero.
\emph{Energy Technology List: Retrofit Measures for Open Fronted Refrigerated Display Cabinets}.
Definitions and product-temperature performance requirements.
\url{https://etl.energysecurity.gov.uk/products/refrigeration-equipment/retrofit-measures-open-fronted-refrigerated-display-cabinets}.
Accessed November 7, 2026.

\bibitem{energystar}
ENERGY STAR. \emph{Ceiling Fan Installation and Usage Tips}.
\url{https://www.energystar.gov/products/ceiling_fans/installation-and-usage-tips}.
Accessed November 7, 2026.

\bibitem{evap}
ASHRAE. \emph{2023 ASHRAE Handbook---HVAC Applications},
Chapter 53, \emph{Evaporative Cooling}.
\url{https://handbook.ashrae.org/Handbooks/A23/SI/A23_Ch53/a23_ch53_si.aspx}.
Accessed November 7, 2026.

\bibitem{problem}
COMAP. \emph{2026 HiMCM Problem B: Air Conditioning the Environment: Modeling Energy Waste in Public Cooling}.
User-supplied problem statement, Requirement 2(a), p.~2.

\bibitem{evapequip}
ASHRAE. \emph{2020 ASHRAE Handbook---HVAC Systems and Equipment},
Chapter 41, \emph{Evaporative Air-Cooling Equipment}, Section 1.
\url{https://handbook.ashrae.org/handbooks/S20/SI/s20_ch41/s20_ch41_si.aspx}.
Accessed November 7, 2026.

\bibitem{fanguide}
Center for the Built Environment, University of California, Berkeley.
\emph{Fans for Cooling People Guidebook: Elevated Air Speed and Thermal Comfort}.
\url{https://cbe-berkeley.gitbook.io/fans-guidebook/full-guidebook/elevated-air-speed-and-thermal-comfort}.
Accessed November 7, 2026.

\bibitem{thermalcomfort}
ASHRAE. \emph{2017 ASHRAE Handbook---Fundamentals},
Chapter 9, \emph{Thermal Comfort}, Section 10, Environmental Indices.
\url{https://handbook.ashrae.org/Handbooks/f17/IP/f17_ch09/f17_ch09_ip.aspx}.
Accessed November 7, 2026.

\bibitem{pythermal}
Tartarini, F., and Schiavon, S. (2020).
pythermalcomfort: A Python package for thermal comfort research.
\emph{SoftwareX}, 12, 100578.
\url{https://doi.org/10.1016/j.softx.2020.100578}.
SET API and applicability limits:
\url{https://pythermalcomfort.readthedocs.io/en/stable/documentation/models.html}.
Accessed November 7, 2026.

\bibitem{psychrolib}
Meyer, D., and Thevenard, D. (2019).
PsychroLib: a library of psychrometric functions to calculate thermodynamic properties of air.
\emph{Journal of Open Source Software}, 4(33), 1137.
\url{https://doi.org/10.21105/joss.01137}.
SI psychrometric API:
\url{https://psychrometrics.github.io/psychrolib/api_docs.html}.
Accessed November 7, 2026.

\bibitem{nycshelter}
New York City Department of Transportation.
\emph{Street Design Manual: Bus Stop Shelter}, Description (regular-size shelter).
\url{https://www.nycstreetdesign.info/furniture/bus-stop-shelter}.
Accessed November 7, 2026.

\bibitem{hillphoenix}
Hillphoenix. \emph{CF-R/CF-FS-R/CF-S-R Multi-Deck Merchandiser: Installation \& Operations Manual}, April 2021, Technical Reference, printed p.~4 (PDF p.~8).
\url{https://www.hillphoenix.com/wp-content/uploads/2021/06/CF-S-R-display-case-manual.pdf}.
Product classification: \url{https://www.hillphoenix.com/products/cf}.
Accessed November 7, 2026.

\bibitem{epzone}
EnergyPlus Development Team. \emph{Engineering Reference},
``Basis for the Zone and Air System Integration'' and ``Moisture Predictor-Corrector.''
\href{https://energyplus.readthedocs.io/en/latest/guides/engineering-reference/2.1-basis-for-the-zone-and-air-system-integration.html}{Online source};
\href{https://energyplus.readthedocs.io/en/latest/guides/engineering-reference/2.4-moisture-predictor-corrector.html}{Online source}.
Accessed November 7, 2026.

\bibitem{cruz2026}
Cruz-Salas, M.V., Puerto, E., Rojas, J., Barrios, G., Huelsz, G., and Ram\'irez, G. (2026).
Mist evaporative cooling system with natural ventilation in a semi-outdoor environment: Theoretical models, experiments, and thermal simulations.
\emph{Energy and Buildings}, 366, 117700.
\url{https://doi.org/10.1016/j.enbuild.2026.117700}.

\bibitem{ventilation}
ASHRAE. \emph{2025 ASHRAE Handbook---Fundamentals}, Chapter 16,
\emph{Ventilation and Infiltration}, Section 10, Natural Ventilation.
\href{https://handbook.ashrae.org/Handbooks/F25/SI/F25_Ch16/F25_Ch16_si.aspx}{Online source}.
Accessed November 7, 2026.

\bibitem{canberra-rsp}
RSP. \emph{Canberra Station}. Project description and site area.
\href{https://rsp.design/project/canberrastation/}{Online source}. Accessed November 7, 2026.

\bibitem{canberra-lta}
Land Transport Authority. \emph{Stations, stations on the tracks, which is the greenest of them all?}
Canberra station: natural ventilation, planting, glazing and louvers.
\href{https://www.lta.gov.sg/content/ltagov/en/who_we_are/statistics_and_publications/Connect/greenmrtstations.html}{Online source}.
Accessed November 7, 2026.

\bibitem{canberra-dimensions}
\emph{The Straits Times}. \emph{Safety first as work starts on new Canberra MRT station on `live' North-South Line}, March 26, 2016.
Source identified in the user-supplied schematic for approximate structural dimensions.
\href{https://www.straitstimes.com/singapore/safety-first-as-work-starts-on-new-canberra-mrt-station-on-live-north-south-line}{Online source}.

\bibitem{singapore-climate}
Meteorological Service Singapore. \emph{Climate of Singapore}. Temperature and humidity climatology.
\href{https://www.weather.gov.sg/climate-climate-of-singapore}{Online source}. Accessed November 7, 2026.

\bibitem{gemini}
Google Gemini. \emph{AI-generated analysis of the 2026 HiMCM Problem B},
user-supplied PDF conversation excerpt, November 2026.
Consulted for process-oriented exposition; the proposed efficiency formula was
reviewed but not adopted. Exact model identifier was not available in the supplied excerpt.

\bibitem{codex}
OpenAI. \emph{Codex}, AI-assisted formulation, reference checking,
English drafting, and LaTeX preparation, November 2026.
Exact deployed model identifier to be recorded from the session interface.
See the preliminary AI-use record following this draft.

\bibitem{aipolicy}
COMAP. \emph{Use of AI Tools in COMAP Contests}, v102025.
\url{https://www.contest.comap.com/undergraduate/contests/mcm/flyer/Contest_AI_Policy.pdf}.
Accessed November 7, 2026.
\end{thebibliography}

% Add Requirement 3, complete Requirement 4 calibration, and add Requirement 5 before References.
% Replace the working summary after completing the full analysis.
\clearpage
\phantomsection
\addcontentsline{toc}{section}{Report on Use of AI Tools}
\section*{Report on Use of AI Tools}
\noindent\textbf{Preliminary record for this partial draft.}

OpenAI Codex assisted with translation, model formulation, source checking, Python calculations, English drafting, figures and LaTeX preparation for Requirements 1 and 2 during November 2026. Record the exact deployed model identifier from the session interface. The following is a usage summary, not a transcript.

For Requirement 1, the user requested justification of cooling-efficiency quantities and their transfer from food storage to public cooling. Codex organized energy, service quantity, quality and duration, reviewed sources and drafted the explanation. A user-supplied Google Gemini analysis informed the process-oriented organization; its exact model identifier was unavailable.

At the user's request, the 2025 Problem B papers of Teams 16442, 16738 and 16739 were read as structural examples. Their estimates, weighting methods and wording were not adopted. The user also supplied freezer, shelter and Canberra schematics. Codex incorporated them as separate PDFs, checked available dimensional sources and distinguished architectural evidence from modeled occupied area and operating assumptions.

For Requirement 2(a), Codex implemented a three-zone bus scenario with psychrometric calculations and published SET, generated charts and CSV outputs, and developed sensitivity and accounting checks. For Requirement 2(b), Codex implemented dynamic background heat/moisture and surface balances, variable counts, allocation and coverage, comparing candidates against a matched no-device baseline.

After the user supplied Cruz-Salas et al.'s research and requested clearer model inheritance, Codex developed a shared local delivery operator, connected it to background balances, added heat-allocation and airflow sensitivities, tested conservation and limiting cases, reran comparisons, and revised figures and prose. AI review checked equations, units, implementation and figure consistency. The cafeteria study's accuracy was not transferred to Canberra.

At the user's request, Codex subsequently separated cooling-efficiency comparison in Requirements 1--3 from overall-effectiveness decisions in Requirement 4. The physical models and numerical results were retained; efficiency-only charts were generated from the frozen CSVs, and the A3 full-candidate and A36 equal-coverage rankings were distinguished. Residual-discomfort/energy tradeoffs and conditional selection rules were relocated to Requirement 4. Cost and non-thermal experience remain explicit uncalibrated inputs.

All numerical outputs are conditional simulations. No field measurements, real installation inventory or experimental validation were produced; independent student verification has not been asserted. Geometry, airflow, mixing, coverage and occupancy assumptions remain identified in the paper and scenario files.

Before submission, add exact prompts and complete outputs for retained AI contributions, including Gemini, and complete tool/model details required by the contest guidance \cite{aipolicy}. Retain this source and conversation as records. Requirement 3, the cost/experience calibration in Requirement 4, Requirement 5 and the final summary remain unfinished.

\end{document}
